\subsection{Transport and the vertical reference edge}
\label{subsec:peps_witness_transport}

\begin{theorem}[Composition of graph transport]%
    \label{thm:peps_graph_transport_composition}
    \lean{TNLean.PEPS.Edge.map_trans}
    \lean{TNLean.PEPS.Tensor.transport_trans}
    \leanok
    Let $\phi:G\simeq G'$ and $\psi:G'\simeq G''$ be graph
    isomorphisms between graphs with linearly ordered vertex sets.
    For every edge $e$ of $G$, one has
    $\psi(\phi(e))=(\psi\circ\phi)(e)$.
    For every tensor $A$ on $G$, one also has
    $(A^\phi)^\psi=A^{\psi\circ\phi}$.
\end{theorem}
\begin{proof}\leanok
    \uses{def:peps_tensor_transport}
    Both edge maps send the endpoint pair of $e$ to the same pair.
    At a vertex $z$, both transported tensors read the component of $A$
    at $\phi^{-1}(\psi^{-1}(z))$. The incident-edge bijections compose,
    and the identifications of equal bond spaces compose with them.
    Thus the local components agree at every physical and virtual index.
\end{proof}

\begin{theorem}[Pulling a coefficient identity back along a graph isomorphism]%
    \label{thm:peps_edge_witness_pullback}
    \lean{TNLean.PEPS.exists_edgeCoeffIdentityWitness_of_transport}
    \leanok
    Let $\phi$ be a graph isomorphism between two tori whose side
    lengths exceed one, and let $A'=A^\phi$ and $B'=B^\phi$.
    Let $e$ be an edge of the original torus and put $f=\phi(e)$.
    Write $e_-$ and $f_-$ for their first endpoints in the chosen edge ordering, and assume
    $\phi(e_-)=f_-$.
    Suppose $f\in\partial R'$, $f_-\in R'$, all bond dimensions
    of $B'$ are positive, and $B'$ is injective on $R'$ and $(R')^c$.
    Suppose also $D_{A'}(f)=D_{B'}(f)$ and there is an invertible
    matrix $Z$ such that
    \begin{align}
        C_{A'}(R',f,M;\sigma',\tau')
         & =C_{B'}(R',f,Z M Z^{-1};\sigma',\tau')
        \label{eq:peps_witness_image_identity}
    \end{align}
    for every inserted matrix $M$ and every physical configuration
    on $R'$ and its complement, after identifying the equal bond spaces.
    Then $R=\phi^{-1}(R')$ has $e\in\partial R$ and $e_-\in R$;
    all bond dimensions of $B$ are positive, $B$ is injective on $R$
    and $R^c$, and $D_A(e)=D_B(e)$.
    There is an invertible matrix $Y$ such that
    \begin{align}
        C_A(R,e,N;\sigma,\tau)
         & =C_B(R,e,Y N Y^{-1};\sigma,\tau)
        \label{eq:peps_witness_original_identity}
    \end{align}
    for every inserted matrix $N$ and every physical configuration
    on $R$ and its complement, with the equal bond spaces identified.
\end{theorem}
\begin{proof}\leanok
    \uses{thm:peps_regionInsertedCoeff_transport,
        thm:peps_region_injective_identification}
    Transport the configurations in~(\ref{eq:peps_witness_original_identity})
    to $R'$ and apply~(\ref{eq:peps_witness_image_identity}).
    Carrying $Z$ back to the bond of $e$ gives $Y$; this identification
    preserves multiplication and inversion, so it carries conjugation
    by $Z$ to conjugation by $Y$.
    Injectivity of both regions and positivity of every bond dimension
    are preserved by the inverse graph isomorphism.
    Finally, $\phi(e_-)=f_-\in R'$ implies $e_-\in\phi^{-1}(R')$.
\end{proof}

\begin{theorem}[A vertical reference-edge coefficient identity]%
    \label{thm:peps_vertical_arc_window_witness}
    \lean{TNLean.PEPS.exists_verticalEdgeCoeffIdentityWitness_of_normalArcWindows}
    \leanok
    Let $A,B$ be translation-invariant tensors on the $w\times h$
    torus, with positive bond dimensions and the same contracted state.
    Suppose every cyclic $L\times K$ window is injective for each tensor,
    where $L,K>0$, $2L+1\leq w$ and $2K+1\leq h$.
    Then there exist a vertical edge $e$, a region $R$, and an invertible
    matrix $Y$ such that $D_A(e)=D_B(e)$, $e\in\partial R$, $e_-\in R$,
    and $B$ is injective on $R$ and $R^c$.
    They satisfy~(\ref{eq:peps_witness_original_identity}) for every
    inserted matrix and every physical configuration on the two regions.
    This is the vertical reference construction in the two-dimensional
    argument of~\cite{Molnar2018NormalPEPS}.
\end{theorem}
\begin{proof}\leanok
    \uses{thm:peps_horizontal_arc_window_witness,
        thm:peps_torus_coordinate_swap_windows, thm:peps_sameState_transport,
        thm:peps_edge_witness_pullback}
    Exchange the coordinates and apply the horizontal construction
    with window lengths $K,L$, at placement $(1,0)$.
    The size bounds give $1+2K\leq h$ and $2L-1\leq w$,
    so this placement satisfies the required inequalities.
    The distinguished right edge is based at $(K,L-1)$ in the
    exchanged torus. Its inverse image is the upward edge based at
    $(L-1,K)$ in the original torus.
    Since $K+1<h$, neither edge crosses a coordinate seam, and
    coordinate exchange preserves the first endpoint in the chosen
    edge ordering. Pull the horizontal identity and its
    injective regions back by Theorem~\ref{thm:peps_edge_witness_pullback}.
\end{proof}
